% ===================================================================
%  Dodatek Z: Status epistemiczny parametru alpha_K
%              i mapa legacy <-> phi-FP
%  TGP v1 -- Sesja v47b (2026-04-12)
% ===================================================================
\section{Status epistemiczny parametru $\alpha_K$ i~mapa
  legacy$\leftrightarrow\varphi$-FP}%
\label{app:Z-alphaK}

Niniejszy dodatek porz\k{a}dkuje status parametru
$\alpha_K \approx 8{,}56$ po sesjach v45--v47b.
Jego celem nie jest wprowadzenie nowej derywacji numerycznej,
lecz \textbf{rozstrzygni\k{e}cie epistemiczne}, czy
$\alpha_K$ pozostaje parametrem fundamentalnym sektora leptonowego,
czy te\.z ma status parametru \emph{legacy},
zast\k{a}pionego przez kanoniczny j\k{e}zyk
$g_0^*$ i~self-consistent $\varphi$-fixed point.

\medskip
\noindent\fbox{\parbox{\linewidth}{%
\textbf{Nota kanoniczna (2026-04-12).}
Od sesji v47b kanoniczny opis sektora leptonowego TGP u\.zywa
Formulacji~A ($K(g)=g^4$) oraz parametr\'ow
$g_0^e$, $g_0^\mu$, $g_0^\tau$, $g_0^*$.
Parametr $\alpha_K^{\rm old}$ zachowujemy jako element j\k{e}zyka
historycznego, przydatny do translacji starszych wynik\'ow
WKB/bifurkacyjnych na j\k{e}zyk kanoniczny.}}
\medskip

% -------------------------------------------------------------------
\subsection{Dwa wsp\'o\l{}istniej\k{a}ce j\k{e}zyki opisu}
\label{app:Z-two-languages}
% -------------------------------------------------------------------

\begin{definition}[J\k{e}zyk legacy]\label{def:Z-legacy}
Przez \textbf{j\k{e}zyk legacy} rozumiemy opis sektora leptonowego,
w~kt\'orym podstawowymi obiektami s\k{a}:
\begin{equation}
  \alpha_K^{\rm old}, \qquad a_\Gamma, \qquad
  r_{21} \leftrightarrow K_2^*/K_1^*,
\end{equation}
oraz powi\k{a}zane z~nimi dawne skany bifurkacyjne i~energetyczne.
W~tym j\k{e}zyku warto\'s\'c
$\alpha_K^{\rm old} \approx 8{,}5616$
by\l{}a traktowana jako parametr efektywny starszej formulacji ODE.
\end{definition}

\begin{definition}[J\k{e}zyk kanoniczny]\label{def:Z-canonical}
Przez \textbf{j\k{e}zyk kanoniczny} rozumiemy opis,
w~kt\'orym podstawowymi obiektami s\k{a}:
\begin{equation}
  g_0^e, \qquad g_0^\mu, \qquad g_0^\tau, \qquad g_0^*,
\end{equation}
oraz twierdzenie self-consistent $\varphi$-fixed point
(tw.~\ref{thm:J2-FP}) wraz z~relacjami dla
$A_{\rm tail}$, $Q_K$ i~widma leptonowego.
\end{definition}

\begin{proposition}[Status relacyjny obu j\k{e}zyk\'ow]
\label{prop:Z-language-status}
Po sesjach v45--v47b:
\begin{enumerate}[nosep]
  \item j\k{e}zyk kanoniczny jest \textbf{silniejszy epistemicznie},
    poniewa\.z koduje sektor leptonowy bez bezpo\'sredniego
    odwo\l{}ania do $\alpha_K^{\rm old}$;
  \item j\k{e}zyk legacy pozostaje \textbf{u\.zyteczny translacyjnie},
    ale nie jest ju\.z minimalnym j\k{e}zykiem opisu;
  \item parametr $\alpha_K^{\rm old}$ ma obecnie status
    \emph{legacy / recovered candidate}, nie za\'s domkni\k{e}tego
    parametru fundamentalnego.
\end{enumerate}
\end{proposition}

\begin{remark}
Wniosek~\ref{prop:Z-language-status} jest sp\'ojny z~uwag\k{a}
\ref{rem:alphaK-legacy-status} w~sek.~08 oraz z~supersesj\k{a}
\'scie\.zki S2 w~Dod.~W (uw.~\ref{rem:W-S2-superseded}).
\end{remark}

% -------------------------------------------------------------------
\subsection{Rozbicie OP-3 na dwa podproblemy}
\label{app:Z-OP3-split}
% -------------------------------------------------------------------

\begin{definition}[Podproblemy OP-3]\label{def:Z-OP3-split}
Problem~\ref{prob:OP-3} rozdzielamy na dwa logicznie roz\l{}\k{a}czne
podproblemy:
\begin{align}
  \textbf{OP-3A}:&\quad
    \text{czy $\alpha_K^{\rm old}$ wynika z~Warstwy~I--II
    bez u\.zycia $r_{21}$ jako wej\'scia?}$, \\
  \textbf{OP-3B}:&\quad
    \text{czy istnieje jawna mapa }
    \alpha_K^{\rm old}\leftrightarrow(g_0^*, g_0^e, \varphi\text{-FP})?
\end{align}
\end{definition}

\begin{remark}[Interpretacja]
OP-3A jest problemem \textbf{silnego wyprowadzenia}
starego parametru.
OP-3B jest problemem \textbf{supersesji formalnej}:
czy dawny parametr mo\.zna odtworzy\'c jako pochodny
obiekt nowego j\k{e}zyka.
\end{remark}

% -------------------------------------------------------------------
\subsection{Bilans dawnych \'scie\.zek selekcji $\alpha_K$}
\label{app:Z-old-paths}
% -------------------------------------------------------------------

\begin{proposition}[Bilans \'scie\.zek legacy]
\label{prop:Z-old-paths-closed}
Dla dawnych dr\'og selekcji $\alpha_K^{\rm old}$ zachodzi:
\begin{enumerate}[nosep]
  \item \textbf{argmin $E^*$}: zamkni\k{e}te negatywnie
    (p96; brak fizycznego minimum przy $r_{\max}\to\infty$),
  \item \textbf{$a_c(\alpha_K)=a_\Gamma$}: zamkni\k{e}te negatywnie
    (p101; brak zgodno\'sci z~punktem fizycznym),
  \item \textbf{$\alpha_K/\alpha_{\rm SN}\sim n_s$}:
    zamkni\k{e}te negatywnie (artefakt p94--p95),
  \item \textbf{argmax $K_2^*/K_1^*$}: zamkni\k{e}te negatywnie,
  \item \textbf{warunki energetyczne}
    $E_{\rm kin}^B/E_{\rm tot}^B$:
    zamkni\k{e}te negatywnie jako selektor fizycznego $\alpha_K$.
\end{enumerate}
W~szczeg\'olno\'sci \'zadna z~tych dr\'og nie daje
domkni\k{e}cia OP-3A.
\end{proposition}

\begin{proof}[Uzasadnienie]
Punkty (1)--(5) s\k{a} bezpo\'srednim bilansem statusu zapisanym
w~sek.~08 (problem~\ref{prob:OP-3}),
\textsc{roadmap} v3 oraz planach analitycznych Koide.
Wszystkie te \'scie\.zki zosta\l{}y albo jawnie obalone,
albo zredukowane do nie-fizycznych koincydencji numerycznych.
\end{proof}

\begin{corollary}[Jedyna sensowna droga bezpo\'srednia]
\label{cor:Z-only-direct-path}
Po odrzuceniu dawnych \'scie\.zek jedyn\k{a} sensown\k{a}
bezpo\'sredni\k{a} drog\k{a} dla OP-3A pozostaje
przepisanie problemu w~postaci implicitnej
\begin{equation}
  g_0(K) + \alpha\,g_1(K) = 0,
\end{equation}
oraz pr\'oba identyfikacji $\alpha$ jako funkcji obiekt\'ow
Warstwy~I--II, a~nie jako parametru skanu dopasowanego do~$r_{21}$.
\end{corollary}

% -------------------------------------------------------------------
\subsection{Status OP-3B: supersesja przez $g_0^*$ i~$\varphi$-FP}
\label{app:Z-OP3B}
% -------------------------------------------------------------------

\begin{proposition}[Cz\k{e}\'sciowe zamkni\k{e}cie OP-3B]
\label{prop:Z-OP3B-partial}
Twierdzenie~\ref{thm:J2-FP} oraz supersesja \'scie\.zki~S2
w~Dod.~W zamykaj\k{a} \textbf{cz\k{e}\'s\'c} problemu OP-3B:
\begin{enumerate}[nosep]
  \item sektor leptonowy posiada ju\.z kanoniczny parametr
    $g_0^*$, kt\'ory zast\k{e}puje dawn\k{a} rol\k{e}
    $\alpha_K^{\rm old}$ jako centralnego obiektu opisu;
  \item relacja $\varphi$-FP odtwarza $r_{21}$ z~dok\l{}adno\'sci\k{a}
    $0{,}0001\%$ bez potrzeby bezpo\'sredniego u\.zycia
    $\alpha_K^{\rm old}$;
  \item kandydat S2c w~Dod.~W pokazuje, \.ze
    $\Phi_0$ mo\.ze by\'c wyznaczane z~$g_0^*$
    bez odwo\l{}ania do starej \'scie\.zki S2.
\end{enumerate}
\end{proposition}

\begin{remark}[Co jeszcze pozostaje otwarte]
OP-3B nie jest jeszcze zamkni\k{e}ty w~pe\l{}ni.
Brakuje nadal:
\begin{enumerate}[nosep]
  \item jawnej formu\l{}y translacyjnej
    $\alpha_K^{\rm old}\leftrightarrow(g_0^*,g_0^e,\varphi\text{-FP})$,
  \item precyzyjnego wskazania, kt\'ore obserwable
    s\k{a} niezmiennicze przy przej\'sciu legacy $\to$ canonical,
  \item zamkni\k{e}tego mostu
    substrat $\to g_0^e$ w~jezyku kanonicznym.
\end{enumerate}
\end{remark}

% -------------------------------------------------------------------
\subsection{Klasyfikacja statusu obiekt\'ow sektora leptonowego}
\label{app:Z-classification}
% -------------------------------------------------------------------

\begin{table}[h]
\centering
\caption{Status epistemiczny obiekt\'ow sektora leptonowego po sesjach v47b}
\label{tab:Z-status}
\begin{tabular}{@{}llll@{}}
\toprule
\textbf{Obiekt} & \textbf{J\k{e}zyk} & \textbf{Status} & \textbf{Komentarz} \\
\midrule
$\alpha_K^{\rm old}$ & legacy & Legacy / Artefakt koordynatowy & OP-3 zamkni\k{e}ty (tw.~\ref{thm:Z-alpha-degeneracy}) \\
$a_\Gamma$ & wsp\'olny & Input / Hipoteza wzmocniona & DESI + Dod.~W \\
$g_0^*$ & kanoniczny & Derived [NUM] & tw.~\ref{thm:J2-FP} \\
$g_0^e$ & kanoniczny & Recovered / Derived & zale\.zny od formulacji \\
$g_0^\mu = \varphi g_0^*$ & kanoniczny & Derived [NUM] & tw.~\ref{thm:J2-FP} \\
$r_{21}$ & obserwabla & Recovered / test & dok\l{}adno\'s\'c $0{,}0001\%$ \\
$Q_K = 3/2$ & kanoniczny & Strukturalny warunek / program domkni\k{e}cia & Dod.~T--T4 \\
\bottomrule
\end{tabular}
\end{table}

% -------------------------------------------------------------------
\subsection{Kryterium domkni\k{e}cia luki}
\label{app:Z-closure}
% -------------------------------------------------------------------

\begin{definition}[Dwa tryby domkni\k{e}cia]\label{def:Z-closure}
Luka zwi\k{a}zana z~$\alpha_K$ mo\.ze zosta\'c uznana za domkni\k{e}t\k{a}
na dwa r\'o\.zne sposoby:
\begin{enumerate}[nosep]
  \item \textbf{Domkni\k{e}cie silne (OP-3A).}
    Istnieje jawny \l{}a\'ncuch
    \[
      \text{Warstwa I--II}
      \to \text{efektywne r\'ownanie solitonu}
      \to \alpha_K^{\rm old}
      \to r_{21},
    \]
    bez u\.zycia $r_{21}$ jako danych wej\'sciowych.
  \item \textbf{Domkni\k{e}cie przez supersesj\k{e} (OP-3B).}
    Parametr $\alpha_K^{\rm old}$ staje si\k{e} obiektem
    \emph{legacy/recovered}, a~kanoniczny most biegnie przez
    \[
      \text{substrat}
      \to g_0^e
      \to \varphi\text{-FP}
      \to (r_{21}, Q_K, m_\tau).
    \]
\end{enumerate}
\end{definition}

% -------------------------------------------------------------------
\subsection{Degeneracja $\alpha$ w~Formie~A: zamkni\k{e}cie OP-3}
\label{app:Z-degeneracy}
% -------------------------------------------------------------------

\begin{theorem}[Degeneracja $\alpha$ -- $\alpha$ jest koordynatowe]
\label{thm:Z-alpha-degeneracy}
W~kanonicznej Formie~A r\'ownania solitonu:
\begin{equation}\label{eq:Z-formA}
  g'' + \frac{d-1}{r}\,g' + \frac{\alpha}{g}\,{g'}^2 + g^2(1-g) = 0,
\end{equation}
obserwable sektora leptonowego s\k{a} \textbf{niezmiennicze pod $\alpha$}:
\begin{align}
  r_{21}(\alpha) &= \mathrm{const} = 206{,}77, \label{eq:Z-r21-inv} \\
  r_{31}(\alpha) &= \mathrm{const} = 3477{,}44, \label{eq:Z-r31-inv} \\
  Q_K(\alpha)    &= \mathrm{const} = \tfrac{3}{2}, \label{eq:Z-QK-inv}
\end{align}
przy kalibracji $\varphi$-FP ($g_0^\mu = \varphi\,g_0^e$)
i~warunku Koide $Q_K = 3/2$.

\medskip\noindent
Jedynym efektem zmiany $\alpha$ jest przesuni\k{e}cie
parametru strza\l{}u $g_0^e(\alpha)$, kt\'ory jest koordynat\k{a}
wewn\k{e}trzn\k{a} ODE, nie obserwablem.
Mapa $\alpha \mapsto g_0^e$ jest monotonicznie malej\k{a}ca
i~odwracalna (diffeomorfizm przestrzeni kalibracji).
\end{theorem}

\begin{proof}[Uzasadnienie numeryczne]
Skrypt \texttt{tgp\_bridge\_op3\_alpha\_degeneracy.py} (v47b)
wykonuje pe\l{}ny skan $\alpha \in [0{,}5;\;3{,}5]$ (7~punkt\'ow),
kalibruj\k{a}c $g_0^e$ metod\k{a} Brenta dla ka\.zdego $\alpha$.
Wyniki (7/7~PASS):
\begin{center}
\begin{tabular}{@{}ccccc@{}}
\toprule
$\alpha$ & $g_0^e$ & $g_0^\mu$ & $r_{31}$ & $Q_K$ \\
\midrule
0{,}50 & 0{,}8967 & 1{,}4509 & 3477{,}44 & 1{,}500000 \\
1{,}00 & 0{,}8879 & 1{,}4367 & 3477{,}44 & 1{,}500000 \\
1{,}50 & 0{,}8783 & 1{,}4210 & 3477{,}44 & 1{,}500000 \\
2{,}00 & 0{,}8677 & 1{,}4040 & 3477{,}44 & 1{,}500000 \\
2{,}50 & 0{,}8562 & 1{,}3854 & 3477{,}44 & 1{,}500000 \\
3{,}00 & 0{,}8436 & 1{,}3650 & 3477{,}44 & 1{,}500000 \\
\bottomrule
\end{tabular}
\end{center}
Spread obserwabli: $\Delta r_{21} = 3 \times 10^{-5}$,
$\Delta r_{31} = 5 \times 10^{-4}$,
$\Delta Q_K = 6 \times 10^{-9}$
--- wszystkie na poziomie b\l{}\k{e}du numerycznego solvera.
\end{proof}

\begin{corollary}[Zamkni\k{e}cie OP-3]
\label{cor:Z-OP3-closed}
Degeneracja tw.~\ref{thm:Z-alpha-degeneracy} zamyka oba
podproblemy OP-3:
\begin{enumerate}[nosep]
  \item \textbf{OP-3A jest bezprzedmiotowy}: $\alpha_K^{\rm old}$
    nie jest parametrem fizycznym lecz artefaktem koordynatowym
    starej Formy~B. Nie ma potrzeby jego wyprowadzania z~Warstwy~I--II.
  \item \textbf{OP-3B jest zamkni\k{e}ty}: jawna mapa
    $\alpha \mapsto (g_0^e(\alpha),\; g_0^\mu = \varphi\,g_0^e,\;
    g_0^\tau\text{ z Koide})$
    daje pe\l{}n\k{a} translacj\k{e} legacy$\leftrightarrow$canonical.
\end{enumerate}
\end{corollary}

\begin{remark}[Predykcja $r_{31}$]
\label{rem:Z-r31-prediction}
Warto\'s\'c $r_{31} = 3477{,}44$ jest \textbf{czyst\k{a} predykcj\k{a}}
TGP (odchylenie $0{,}006\%$ od PDG $3477{,}23$).
Po kalibracji $r_{21}$ sektor leptonowy ma
$N_{\rm free} = 0$ wolnych parametr\'ow.
\end{remark}

\begin{theorem}[Status na dzie\'n 2026-04-12 -- po degeneracji]
\label{thm:Z-main-status}
Na dzie\'n 2026-04-12, po potwierdzeniu degeneracji $\alpha$:
\begin{enumerate}[nosep]
  \item OP-3A jest \textbf{bezprzedmiotowy}
    ($\alpha$ jest koordynatowe);
  \item OP-3B jest \textbf{zamkni\k{e}ty}
    (jawna mapa translacyjna, tw.~\ref{thm:Z-alpha-degeneracy});
  \item $\alpha_K^{\rm old} \approx 8{,}56$ ma status
    \emph{legacy / artefakt koordynatowy};
  \item kanoniczny opis sektora leptonowego biegnie przez
    $g_0^*$, $\varphi$-FP i~warunek Koide $Q_K = 3/2$;
  \item sektor leptonowy TGP ma $N_{\rm free} = 0$
    po kalibracji $r_{21}$, a~$r_{31}$ jest czyst\k{a} predykcj\k{a}.
\end{enumerate}
\end{theorem}

\begin{remark}[Rekomendacja programowa -- aktualizacja]
Po zamkni\k{e}ciu OP-3 najwa\.zniejsze otwarte problemy to:
\begin{enumerate}[nosep]
  \item domkni\k{e}cie mostu substrat $\to g_0^e$,
  \item pe\l{}na analityczna derywacja $\varphi$-FP
    z~pierwszych zasad (Warstwa~I--II),
  \item dalsze testy predykcji $r_{31}$ i~$Q_K$
    wzgl\k{e}dem przysz\l{}ych danych eksperymentalnych.
\end{enumerate}
\end{remark}
